\chapter{Pruebas}\label{pruebas}

\section{Introducci\'on}

La fase de pruebas tiene como objetivo validar el sistema desde una doble perspectiva:
\textbf{correctitud funcional} y \textbf{calidad de implementaci\'on}.
En un proyecto con arquitectura desacoplada (frontend y backend),
la validaci\'on debe cubrir tanto el comportamiento interno de los m\'odulos
como los flujos de usuario de extremo a extremo.

Por este motivo se adopta una estrategia de pruebas en capas,
complementada con an\'alisis de mutaci\'on y ejecuci\'on automatizada en CI/CD.

\section{Objetivos de verificaci\'on}

Los objetivos principales de esta fase son:

\begin{itemize}
	\item verificar que los casos de uso prioritarios funcionan de forma estable;
	\item prevenir regresiones al introducir cambios en l\'ogica de negocio;
	\item garantizar que las restricciones de seguridad y autorizaci\'on se respetan;
	\item medir cobertura y fortaleza de tests de forma objetiva;
	\item integrar la validaci\'on en el ciclo continuo de construcci\'on y despliegue.
\end{itemize}

\section{Estrategia de pruebas en capas}

La estrategia sigue una pir\'amide de testing con cuatro niveles:

\begin{enumerate}
	\item \textbf{Pruebas unitarias}: validan l\'ogica aislada de servicios, utilidades y componentes.
	\item \textbf{Pruebas de integraci\'on}: validan interacci\'on entre capas y contratos de API.
	\item \textbf{Pruebas E2E}: validan flujos de usuario en navegador.
	\item \textbf{Pruebas de mutaci\'on}: eval\'uan la capacidad real de los tests para detectar defectos.
\end{enumerate}

La Tabla~\ref{tab:matriz-pruebas} resume las tecnolog\'ias aplicadas.

\begin{table}[ht]
\centering
\begin{tabular}{|P{2.9cm}|P{3.6cm}|P{5.4cm}|}
\hline
\textbf{Nivel} & \textbf{Herramientas} & \textbf{Finalidad principal} \\
\hline
Unitarias frontend & Vitest + Testing Library & Validaci\'on de componentes, contextos, servicios y utilidades del cliente. \\
\hline
Unitarias e integraci\'on backend & JUnit + Spring Test + MockMvc & Verificar servicios, controladores, seguridad y persistencia. \\
\hline
E2E frontend & Playwright & Simular escenarios reales de uso en navegador con flujos completos. \\
\hline
Mutaci\'on frontend & Stryker & Medir fortaleza de tests frente a mutantes en servicios y utilidades. \\
\hline
Mutaci\'on backend & PIT & Medir capacidad de detecci\'on de fallos en servicios/controladores objetivo. \\
\hline
\end{tabular}
\caption{Matriz de niveles y herramientas de prueba}
\label{tab:matriz-pruebas}
\end{table}

\section{Inventario cuantitativo de pruebas realizadas}

Para documentar de forma expl\'icita el alcance de verificaci\'on,
se ha consolidado el recuento de pruebas por capa a partir de los reportes
generados por Vitest, Playwright y Surefire.

El n\'umero total de pruebas autom\'aticas ejecutadas en el \textit{corte} actual es:

\[
N_{\text{total}} = N_{\text{frontend unitarias}} + N_{\text{frontend E2E}} + N_{\text{backend}} = 137 + 3 + 816 = 956
\]

Desglose principal:

\begin{itemize}
\item frontend unitarias: 137 casos de prueba en 21 ficheros;
\item frontend E2E: 3 escenarios end-to-end;
\item backend: 816 casos de prueba agregados en 56 suites Surefire.
\end{itemize}

\section{Pruebas en frontend}

\subsection{Pruebas unitarias y cobertura}

El frontend utiliza Vitest como runner principal,
con entorno \texttt{happy-dom} y configuraci\'on de cobertura por V8.
La suite incluye pruebas sobre:

\begin{itemize}
	\item componentes (por ejemplo, navegaci\'on y constructores de estructura ZIP),
	\item contextos de autenticaci\'on y tema,
	\item p\'aginas clave (login, dashboard, evaluaciones, perfil),
	\item servicios de API (auth, cursos, actividades, entregables, entregas, feedback),
	\item utilidades de parsing/validaci\'on de estructura ZIP.
\end{itemize}

Se han definido umbrales m\'inimos de cobertura en la configuraci\'on de Vitest:
\begin{itemize}
	\item \texttt{lines} \(\ge 95\%\),
	\item \texttt{functions} \(\ge 95\%\),
	\item \texttt{statements} \(\ge 95\%\),
	\item \texttt{branches} \(\ge 85\%\).
\end{itemize}

\subsection{M\'odulos testeados en frontend}

Adem\'as de la cobertura porcentual, se ha realizado un inventario por m\'odulo
de la suite unitaria para dejar evidencia de qu\'e zonas funcionales reciben
mayor esfuerzo de verificaci\'on.

\begin{table}[ht]
\centering
\begin{tabular}{|P{5.6cm}|P{2.2cm}|}
\hline
\textbf{M\'odulo frontend} & \textbf{Casos} \\
\hline
Servicios (\texttt{src/services}) & 51 \\
\hline
Componentes (src/components) & 30 \\
\hline
Utilidades (\texttt{src/utils}) & 22 \\
\hline
P\'aginas (\texttt{src/pages}) & 15 \\
\hline
Contextos (\texttt{src/context}) & 11 \\
\hline
Ra\'iz de aplicaci\'on (\texttt{src/}) & 8 \\
\hline
\textbf{Total frontend unitarias} & \textbf{137} \\
\hline
\end{tabular}
\caption{Distribuci\'on de casos unitarios por m\'odulo en frontend}
\label{tab:frontend-modulos-testeados}
\end{table}

\subsection{Pruebas E2E}

Las pruebas end-to-end se ejecutan con Playwright sobre Chromium,
y validan flujos de interacci\'on completos desde la interfaz.

El enfoque de ejecuci\'on adoptado en este proyecto es \textbf{E2E mockeada}:
el navegador se levanta contra el frontend real y se simulan endpoints API
en los escenarios definidos. Esto reduce inestabilidad ambiental y facilita
reproducibilidad en CI.

Se validan, entre otros, los escenarios:

\begin{itemize}
	\item autenticaci\'on correcta,
	\item autenticaci\'on fallida,
	\item flujo de evaluaciones con filtros y descarga.
\end{itemize}

\subsection{Mutaci\'on en frontend}

Stryker se aplica sobre c\'odigo de producci\'on de servicios y utilidades,
excluyendo archivos de test.
El objetivo es verificar que los tests no solo ejecutan l\'ineas,
sino que detectan alteraciones sem\'anticas en el comportamiento.

La cobertura mediante mutaci\'on en frontend se ha realizado con el siguiente
procedimiento:

\begin{enumerate}
\item selecci\'on del alcance de mutaci\'on en \path{src/services/**/*.ts} y \path{src/utils/**/*.ts};
\item exclusi\'on expl\'icita de tests y archivos de infraestructura (\path{*.test.*}, \path{*.spec.*}, \path{api.ts}, \path{index.ts});
\item ejecuci\'on de \texttt{npm run test:mutation}, generando mutantes y relanzando Vitest para cada lote;
\item an\'alisis de mutantes \textit{survived}/\textit{timeout} para reforzar casos de prueba d\'ebiles.
\end{enumerate}

Como criterio de aceptaci\'on de esta fase, Stryker aplica umbrales
\texttt{high=75}, \texttt{low=60} y \texttt{break=50},
siendo este \'ultimo el que determina fallo de pipeline.

\section{Pruebas en backend}

\subsection{Pruebas unitarias e integraci\'on}

En backend se emplea JUnit junto con Spring Test y MockMvc.
La validaci\'on cubre:

\begin{itemize}
	\item l\'ogica de servicios de dominio,
	\item validaciones y reglas de negocio,
	\item contratos REST de controladores,
	\item seguridad y autorizaci\'on por contexto,
	\item entidades y conversiones DTO.
\end{itemize}

Entre los tests de referencia de mayor carga funcional destaca
\texttt{EntregaServiceTest}, que valida escenarios de entrega,
restricciones y evoluci\'on de estado.

\subsection{M\'odulos testeados en backend}

El backend se ha cuantificado a partir de nodos \texttt{testcase}
en reportes Surefire, lo que permite reflejar correctamente pruebas
parametrizadas y escenarios anidados.

\begin{table}[ht]
\centering
\begin{tabular}{|P{5.6cm}|P{2.2cm}|}
\hline
\textbf{M\'odulo backend} & \textbf{Casos} \\
\hline
Servicios (\texttt{service}) & 513 \\
\hline
Controladores (\texttt{controller}) & 144 \\
\hline
Repositorios (\texttt{repository}) & 47 \\
\hline
DTOs y validaciones (\texttt{dto}) & 40 \\
\hline
Seguridad (\texttt{security}) & 31 \\
\hline
Entidades (\texttt{entity}) & 22 \\
\hline
Manejo de excepciones (exception) & 10 \\
\hline
Configuraci\'on (\texttt{config}) & 8 \\
\hline
Aplicaci\'on base (\texttt{root}) & 1 \\
\hline
\textbf{Total backend} & \textbf{816} \\
\hline
\end{tabular}
\caption{Distribuci\'on de casos de prueba por m\'odulo en backend}
\label{tab:backend-modulos-testeados}
\end{table}

\subsection{Cobertura backend}

La cobertura se genera con JaCoCo durante la fase de test,
produciendo informe HTML en
\texttt{backend/target/site/jacoco/index.html}.

Este reporte se utiliza como soporte de revisi\'on t\'ecnica
para detectar m\'odulos subprobados y orientar ampliaciones de suite.

\subsection{Mutaci\'on en backend}

El perfil Maven \texttt{mutation} ejecuta PIT sobre clases objetivo
de servicios y controladores seleccionados.
La configuraci\'on define umbrales de mutaci\'on y cobertura,
y genera reportes estructurados para su an\'alisis.

En esta implementaci\'on se mutan de forma controlada 10 clases objetivo
(7 de servicio y 3 de controlador). Las clases objetivo son:

\begin{itemize}
\item \texttt{OneDriveService}
\item \texttt{ZipValidationService}
\item \texttt{FeedbackService}
\item \texttt{SecurityContextUserService}
\item \texttt{CloudinaryService}
\item \texttt{OneDriveController}
\item \texttt{MicrosoftOAuthController}
\item \texttt{FeedbackController}
\end{itemize}

El flujo seguido para cobertura de mutaci\'on en backend ha sido:

\begin{enumerate}
\item ejecuci\'on del perfil con \texttt{mvn -Pmutation ... mutationCoverage};
\item generaci\'on de mutantes con operadores de retorno, aritm\'etica e incrementos;
\item relaci\'on mutante--test para identificar qu\'e prueba mata cada defecto sint\'etico;
\item revisi\'on de mutantes supervivientes para crear pruebas adicionales o ajustar aserciones.
\end{enumerate}

El perfil exige adem\'as un m\'inimo de mutaci\'on del 70\%
y de cobertura del 75\%, reforzando que la cobertura no sea
solo de l\'ineas, sino de detecci\'on efectiva de fallos.

\section{Automatizaci\'on en CI/CD}

La validaci\'on se integra en pipelines de GitHub Actions,
lo que garantiza ejecuci\'on repetible por cada cambio relevante.

Workflows principales:

\begin{itemize}
	\item \texttt{frontend-ci.yml}: lint, unitarias con cobertura, E2E y build de frontend.
	\item \texttt{full-stack-ci.yml}: verificaci\'on backend (\texttt{mvn verify}) y frontend completo.
	\item \texttt{mutation-ci.yml}: mutaci\'on backend (PIT) y frontend (Stryker), con subida de artefactos.
\end{itemize}

Esta automatizaci\'on reduce el riesgo de integraci\'on tard\'ia,
y permite detectar regresiones antes de despliegues a entornos superiores.

\section{Resultados verificados}

De acuerdo con el informe de estrategia de testing del proyecto
(corte de verificaci\'on final), se obtienen los siguientes resultados representativos:

\begin{table}[ht]
\centering
\begin{tabular}{|P{4.4cm}|P{7.4cm}|}
\hline
\textbf{M\'etrica} & \textbf{Resultado} \\
\hline
E2E frontend (Playwright) & 3 escenarios ejecutados, 3 en verde. \\
\hline
Cobertura frontend & Statements 97.77\%, Branches 90.68\%, Functions 98.19\%, Lines 98.25\%. \\
\hline
Suite frontend & 137 tests aprobados en 21 archivos durante la ejecuci\'on del comando test:coverage. \\
\hline
Suite backend completa & 816 casos en 56 suites Surefire, 0 fallos y 0 errores. \\
\hline
Total pruebas autom\'aticas & 956 ejecuciones (137 unitarias frontend + 3 E2E + 816 backend). \\
\hline
Mutaci\'on frontend (Stryker) & Score global 98.61\% (por encima del umbral \texttt{break}=50). \\
\hline
Mutaci\'on backend (PIT) & Mutation score 93\%, line coverage 95\%, test strength 96\%. \\
\hline
Test objetivo backend & \texttt{EntregaServiceTest}: 89 tests, 0 fallos, 0 errores. \\
\hline
\end{tabular}
\caption{Resumen de resultados de verificaci\'on}
\label{tab:resultados-pruebas}
\end{table}

Estos valores muestran una madurez de testing consistente con el nivel de exigencia
de un TFG de orientaci\'on profesional.

\section{Criterios de aceptaci\'on}

Como criterio operativo de calidad se establecen los siguientes m\'inimos:

\begin{itemize}
	\item \texttt{lint} sin errores,
	\item suites unitarias/integraci\'on/E2E en verde,
	\item coberturas por encima de umbrales configurados,
	\item m\'etricas de mutaci\'on por encima de umbral de ruptura,
	\item ausencia de regresiones funcionales en flujos cr\'iticos.
\end{itemize}

\section{Limitaciones y riesgos}

Aunque la estrategia es robusta, se identifican l\'ineas de mejora:

\begin{itemize}
	\item ampliar E2E con backend real en entorno de staging para ciertos casos,
	\item extender mutaci\'on a mayor superficie de clases backend,
	\item reforzar pruebas de carga y comportamiento en condiciones extremas.
\end{itemize}

No obstante, las limitaciones actuales no comprometen la validez de la verificaci\'on
para el alcance definido del proyecto.

\section{Conclusiones del cap\'itulo}

La estrategia de pruebas aplicada combina cobertura, profundidad y automatizaci\'on,
aportando evidencia objetiva de calidad funcional y t\'ecnica.

La incorporaci\'on de mutaci\'on como criterio adicional permite evaluar
la eficacia real de la suite, superando una visi\'on meramente cuantitativa
de cobertura.

En conjunto, los resultados respaldan que el sistema dispone de una base
de validaci\'on s\'olida para mantenimiento, evoluci\'on y despliegue continuo.






